
\documentclass[12pt]{article} % use larger type; default would be 10pt
\usepackage[square]{natbib}
\usepackage[utf8]{inputenc} % set input encoding (not needed with XeLaTeX)
\usepackage{amsmath,amssymb}
\DeclareMathOperator{\E}{\mathbb{E}}
\usepackage{color}
\usepackage{pdfpages}
%%% PAGE DIMENSIONS
\usepackage{geometry} % to change the page dimensions
\geometry{a4paper} % or letterpaper (US) or a5paper or....
 \geometry{margin=1in} % for example, change the margins to 2 inches all round
% \geometry{landscape} % set up the page for landscape
%   read geometry.pdf for detailed page layout information
 \newcommand{\eps}{\varepsilon}
\usepackage{graphicx} % support the \includegraphics command and options

% \usepackage[parfill]{parskip} % Activate to begin paragraphs with an empty line rather than an indent

%%% PACKAGES
\usepackage{booktabs} % for much better looking tables
\usepackage{array} % for better arrays (eg matrices) in maths
\usepackage{paralist} % very flexible & customisable lists (eg. enumerate/itemize, etc.)
\usepackage{verbatim} % adds environment for commenting out blocks of text & for better verbatim
\usepackage{subfig} % make it possible to include more than one captioned figure/table in a single float
% These packages are all incorporated in the memoir class to one degree or another...

%%% HEADERS & FOOTERS
\usepackage{fancyhdr} % This should be set AFTER setting up the page geometry
%\pagestyle{fancy} % options: empty , plain , fancy

\usepackage{tikz}

\fancypagestyle{firstpage}
{
\renewcommand{\headrulewidth}{1pt} % customise the layout...
\lhead{STAT 400 (FA26)}\chead{HW 07}\rhead{}
\lfoot{}\cfoot{\thepage}\rfoot{}
}


%%% SECTION TITLE APPEARANCE
\usepackage{sectsty}
\allsectionsfont{\sffamily\mdseries\upshape} % (See the fntguide.pdf for font help)
% (This matches ConTeXt defaults)

%%% ToC (table of contents) APPEARANCE
\usepackage[nottoc,notlof,notlot]{tocbibind} % Put the bibliography in the ToC
\usepackage[titles,subfigure]{tocloft} % Alter the style of the Table of Contents
\renewcommand{\cftsecfont}{\rmfamily\mdseries\upshape}
\renewcommand{\cftsecpagefont}{\rmfamily\mdseries\upshape} % No bold!


\definecolor{Red}{rgb}{1,0,0}
\newcommand{\Red}{\color{Red}}
\definecolor{Blue}{rgb}{0,0,1}
\newcommand{\Blue}{\color{Blue}}


%%%%%%%THEOREM 
\usepackage{thmtools,thm-restate}
\newtheorem{theorem}{Theorem}
\newtheorem{lemma}[theorem]{Lemma}
\newtheorem{proposition}[theorem]{Proposition}
\newtheorem{corollary}[theorem]{Corollary}

\newenvironment{proof}[1][Proof]{\begin{trivlist}
\item[\hskip \labelsep {\bfseries #1}]}{\end{trivlist}}
\newenvironment{definition}[1][Definition]{\begin{trivlist}
\item[\hskip \labelsep {\bfseries #1}]}{\end{trivlist}}
\newenvironment{example}[1][Example]{\begin{trivlist}
\item[\hskip \labelsep {\bfseries #1}]}{\end{trivlist}}
\newenvironment{remark}[1][Remark]{\begin{trivlist}
\item[\hskip \labelsep {\bfseries #1}]}{\end{trivlist}}

\newcommand{\qed}{\nobreak \ifvmode \relax \else
      \ifdim\lastskip<1.5em \hskip-\lastskip
      \hskip1.5em plus0em minus0.5em \fi \nobreak
      \vrule height0.75em width0.5em depth0.25em\fi}
   
     \usepackage{tikz}
  \def\firstcircle{(90:1cm) circle (1.5cm)}
  \def\secondcircle{(210:1cm) circle (1.5cm)}
  \def\thirdcircle{(330:1cm) circle (1.5cm)}
   
   
%\usepackage{nopageno}   
 %%%%%%%%%%%%%%

\linespread{2}


\usepackage{enumitem}



%%% END Article customizations



%%% The "real" document content comes below...






\begin{document}
%\maketitle
\thispagestyle{firstpage}


\begin{enumerate}
\item Suppose that number of accidents at the Monstropolis power plant follows the Poisson process with the average rate of 0.6 accidents per day. 
\begin{enumerate}[label=1.\arabic*]
\item Use Normal approximation to find the probability that at most 200 accidents will occur in a year (365 days). \\ {\Red Let $X$ be the number of accidents that occur in a year. We have $X \sim \text{Poisson}(219)$. By normal approximation, we have $$\begin{aligned}
    P(X \leq 200) &= P\left( Z \leq \frac{200+0.5-219}{\sqrt{219}} \right) \\
    &= 0.1056
\end{aligned}$$}
\item Use Normal approximation to find the probability of less than 190 accident-free days in a year (365 days). \\ {\Red Let $Y$ denote the number of accident-free days in a year. We have $Y \sim \text{Binom}(365, p)$ where $$\begin{aligned}
    p &= P(\text{Poisson}(0.6) = 0) \\
    &= e^{-0.6}
\end{aligned}$$ By normal approximation, we have $$\begin{aligned}
    P(Y < 190) &= P(Y \leq 189) \\
    &= P\left(Z \leq \frac{189+0.5-365 \cdot e^{-0.6}}{\sqrt{365 \cdot e^{-0.6}(1-e^{-0.6})}}\right) \\
    &= 0.1276
    \end{aligned}$$}
    (R: \texttt{p <- exp(-0.6); pnorm(189.5, 365*p, sqrt(365*p*(1-p)))} gives 0.1276)
\end{enumerate}
\newpage
\item The distribution of the baggage weights for passengers using a particular airline has a mean of $21$ lbs and a standard deviation of $5.3$ lbs. 
\begin{enumerate}[label=2.\arabic*]
\item What is the probability that for (a random sample of) 100 passengers, the total luggage weight is less than 2,100 lbs? \\ {\Red $$\begin{aligned}
    P(S < 2100) &= P(\bar{X} < 21) \\
    &= P(Z < 0) \\
    &= 0.5
\end{aligned}$$}
\item What is the probability that for (a random sample of) 100 passengers, the sample mean weight is within 0.5lb of the overall mean (21 lbs)? {\Red $$\begin{aligned}
    P(20.5 < \bar{X} < 21.5) &= P\left(\frac{20.5-21}{5.3/\sqrt{100}} < \frac{\bar{X}-21}{5.3/\sqrt{100}} < \frac{21.5-21}{5.3/\sqrt{100}}\right) \\
    &= P(-0.94 < Z < 0.94) \\
    &= 2(0.8264) - 1 \\
    &= 0.6528
\end{aligned}$$}
(R: \texttt{pnorm(21.5, 21, 0.53) - pnorm(20.5, 21, 0.53)} gives 0.6545)
\end{enumerate}
\newpage
\item Let $\bar{X}$ be the mean of a random sample of size $n=15$ from a distribution with mean $\mu=80$ and variance $\sigma^2=60$.  Use the Chebyshev's inequality to find a lower bound for $P(74<\bar{X}<86)$. \\ {\Red $$\begin{aligned}
    P(74 < \bar{X} < 86) &= P(|\bar{X}-80| < 6) \\
    &\geq 1 - \frac{4}{36} \\
    &= \frac{8}{9}
\end{aligned}$$}
\newpage
\item Let $X_{1}$ and $X_{2}$ be independent Poisson random variables with respective means $\lambda_1=4$ and $\lambda_2=3$. 
\begin{enumerate}[label=4.\arabic*]
\item Find $P(X_1=3,X_2=5)$. \\ {\Red $$\begin{aligned}
    P(X_1 = 3, X_2 = 5) &= P(X_1 = 3) P(X_2 = 5) \\
    &= 0.0197
\end{aligned}$$}
\item Find $P(X_1+X_2=1)$. \\ {\Red Since $X_1$ and $X_2$ are independent, we have $X_1 + X_2 \sim \text{Pois}(7)$. So we have $$\begin{aligned}
    P(X_1 + X_2 = 1) &= 0.0064
\end{aligned}$$}
\end{enumerate}
\newpage
\item Let $X_{1}$ and $X_{2}$ be independent binomial random variables, such that $X_1\sim Bin(n=4, p=0.3)$ and $X_2\sim Bin(n=6, p=0.3)$.
\begin{enumerate}[label=5.\arabic*]
\item Find $P(X_1=2,X_2=4)$. \\ {\Red Since $X_1$ and $X_2$ are independent, we have $$\begin{aligned}
    P(X_1 = 2, X_2 = 4) &= P(X_1 = 2) P(X_2 = 4) \\
    &= 0.0158
\end{aligned}$$}
\item Find $P(X_1+X_2=5)$. \\ {\Red Since $X_1$ and $X_2$ are independent, we have $X_1+X_2 \sim \text{Binom}(10, 0.3)$. So we have $$\begin{aligned}
    P(X_1 + X_2 = 5) &= 0.1029
\end{aligned}$$}
\end{enumerate}
\newpage
\item Let $X_1, \dots, X_n$ be a random sample of size $n$ from a geometric distribution for which $p$ is the probability of success. 

Use the following data to give a point estimate (maximum likelihood estimate) of $p$:
%
\begin{center}
\begin{tabular}{c c c c c} 
4& 25& 17& 16&10\\
7& 6& 20& 11& 1
\end{tabular}
\end{center}
%
\textcolor{red}{The pmf of the geometric distribution is $f(x;p) = (1-p)^{x-1}p$ for $x = 1, 2, \dots$. We start by finding the likelihood $$\begin{aligned}
    \mathcal{L}(p) &= f_{X_1, \dots, X_n}(x_1, \dots, x_n; p) \\
    &= \prod_{i=1}^n (1-p)^{x_i - 1} p \\
    &= p^n (1-p)^{\sum_{i=1}^n x_i - n}
\end{aligned}$$ Taking the log likelihood, we get $$\begin{aligned}
    l(p) &= n \log(p) + \left(\sum_{i=1}^n x_i - n\right) \log(1-p)
\end{aligned}$$ We maximize the log likelihood by taking the first order condition $$\begin{aligned}
    l'(p) &= \frac{n}{p} - \frac{\sum_{i=1}^n x_i - n}{1-p} \\
    &\overset{\mathrm{set}}{=} 0
\end{aligned}$$ Solving for $p$ yields $$\begin{aligned}
    n(1-p) &= p\left(\sum_{i=1}^n x_i - n\right) \\
    n &= p \sum_{i=1}^n x_i \\
    \hat p &= \frac{n}{\sum_{i=1}^n x_i} = \frac{1}{\bar{x}}
\end{aligned}$$ Since $l''(p) = -\frac{n}{p^2} - \frac{\sum x_i - n}{(1-p)^2} < 0$, this is a maximum. From the data, we have $n = 10$ and $\sum_{i=1}^{10} x_i = 117$, so $\bar{x} = 11.7$. Therefore, the maximum likelihood estimate is $$\begin{aligned}
    \hat p = \frac{10}{117} = \frac{1}{11.7} \approx 0.0855
\end{aligned}$$}

\newpage
% \item An urn contains 64 balls, of which $N_1$ are orange and $N_2$ are blue. A random sample of $n=8$ balls is selected from the urn without replacement, and $X$ is equal to the number of orange balls in the sample. This experiment was repeated 30 times (the 8 balls being returned to the urn before each repetition), yielding the following data:
% \begin{center}
% \begin{tabular}{c c c c c} 
% 0& 0& 2& 1& 1\\
% 4& 3& 1& 1& 2 \\
% 2& 2& 1& 0& 2\\
% 2& 1& 2& 4& 2\\
% 0& 1& 1& 1& 3\\
% 0& 2& 3& 0& 1
% \end{tabular}
% \end{center}
% \begin{enumerate}[label=7.\arabic*]
% \item What is the distribution of this experiment's outcome?
% \item Using these data, guess the value of $N_1$ to the nearest whole number using method of moment method. 
% \end{enumerate}

\newpage
\item An urn contains 64 balls, of which $N_1$ are orange and $N_2$ are blue. A random sample of $n=8$ balls is selected from the urn without replacement, and $X$ is equal to the number of orange balls in the sample. 
\begin{enumerate}[label=7.\arabic*]
\item What is the distribution of this experiment's outcome? \\ \textcolor{red}{Hypergeometric.}
\item This experiment was repeated 30 times (the 8 balls being returned to the urn before each repetition of an experiment), yielding the following data:
\begin{center}
\begin{tabular}{c c c c c} 
0& 0& 2& 1& 1\\
4& 3& 1& 1& 2 \\
2& 2& 1& 0& 2\\
2& 1& 2& 4& 2\\
0& 1& 1& 1& 3\\
0& 2& 3& 0& 1
\end{tabular}
\end{center}
Using these data, guess the value of $N_1$ to the nearest whole number using method of moment method. \\ 
\textcolor{red}{By the properties of the hypergeometric distribution, we know that $\mathbb{E}(X) = \frac{8N_1}{64} = \frac{N_1}{8}$. Setting the first sample moment equal to the first population moment, we get $$\begin{aligned}
    \bar{x} = \frac{N_1}{8}
\end{aligned}$$ Solving for $N_1$ yields $$\begin{aligned}
    \hat{N_1} = 8 \bar{x}
\end{aligned}$$ From the data, we know that $\bar{x} = 1.5$. Therefore, the method of moment estimate is $$\begin{aligned}
    \hat N_1 &= 12
\end{aligned}$$}
\end{enumerate}
\newpage
\item Given the following 25 observations from a
gamma distribution$(\alpha, \theta)$, use the method-of-moments estimators to find point estimates of $\alpha$ and $\theta$:

% 6.9 7.3 6.7 6.4 6.3 
% 5.9 7.0 7.1 6.5 7.6 
% 7.2 7.1 6.1 7.3 7.6 
% 7.6 6.7 6.3 5.7 6.7 
% 7.5 5.3 5.4 7.4 6.9
\begin{center}
\begin{tabular}{ccccc}
6.9 & 7.3 & 6.7 & 6.4 & 6.3 \\
5.9 & 7.0 & 7.1 & 6.5 & 7.6 \\
7.2 & 7.1 & 6.1 & 7.3 & 7.6 \\
7.6 & 6.7 & 6.3 & 5.7 & 6.7 \\
7.5 & 5.3 & 5.4 & 7.4 & 6.9 \\
\end{tabular}
\end{center}
\textcolor{red}{By the properties of the Gamma distribution, we have $\mathbb{E}(X) = \alpha \theta$. Setting the first sample moment equal to the first population moment, we have $$\begin{aligned}
    \bar{x} = \alpha \theta 
\end{aligned}$$ Since there are unknown parameters to estimate, we need two equations. So let's also set the second sample moment equal to the second population moment $\mathbb{E}(X^2) = \alpha \theta^2 + \alpha^2 \theta^2$. That gives us $$\begin{aligned}
    \sum_{i=1}^n x_i^2/n = \alpha \theta^2 + \alpha^2 \theta^2
\end{aligned}$$ Solving for $\alpha$ and $\theta$ yields $$\begin{aligned}
    \hat \alpha &= \frac{n \bar{x}^2}{\sum_{i=1}^n x_i^2 - n\bar{x}^2} \\
    \hat \theta &= \frac{\sum_{i=1}^n x_i^2 - n \bar{x}^2}{n \bar{x}}
\end{aligned}$$Plugging the sample in yields $$\begin{aligned}
    \hat \alpha &= 102.50 \\
    \hat \theta &= 0.06576
\end{aligned}$$}

\newpage
\item Find the maximum likelihood estimates for $\theta_1=\mu$
and $\theta_2 = \sigma^2$ if a random sample of size 15 from $N(\mu,\sigma^2)$
yielded the following values:
% 31.5 36.9 33.8 30.1 33.9
% 35.2 29.6 34.4 30.5 34.2
% 31.6 36.7 35.8 34.5 32.7

\begin{center}
    \begin{tabular}{ccccc}
31.5 & 36.9 & 33.8 & 30.1 & 33.9 \\
35.2 & 29.6 & 34.4 & 30.5 & 34.2 \\
31.6 & 36.7 & 35.8 & 34.5 & 32.7 \\
\end{tabular}
\end{center}
\textcolor{red}{We start by finding the likelihood $$\begin{aligned}
    \mathcal{L}(\theta_1, \theta_2) &= f_{X_1, X_2, ..., X_n}(x_1, x_2, ..., x_n; \theta_1, \theta_2) \\
    &= \prod_{i=1}^n f(x_i; \theta_1, \theta_2) \\
    &= \prod_{i=1}^n \frac{1}{\sqrt{2\pi\theta_2}} e^{-\frac{(x_i-\theta_1)^2}{2\theta_2}} \\
    &= \frac{1}{(2\pi)^{n/2} \theta_2^{n/2}} e^{-\frac{\sum_{i=1}^n (x_i-\theta_1)^2}{2\theta_2}}
\end{aligned}$$ Taking the log likelihood, we get $$\begin{aligned}
    l(\theta_1, \theta_2) &= -\frac{n}{2} \log(2\pi) - \frac{n}{2} \log(\theta_2) - \frac{\sum_{i=1}^n(x_i-\theta_1)^2}{2 \theta_2}  
\end{aligned}$$ We maximize the log likelihood by taking the first order conditions $$\begin{aligned}
    \frac{\partial l}{\partial \theta_1} &= \frac{\sum_{i=1}^n x_i - n\theta_1}{\theta_2} \\
    &\overset{\mathrm{set}}{=} 0 \\
    \frac{\partial l}{\partial \theta_2} &= -\frac{n}{2\theta_2} + \frac{\sum_{i=1}^n (x_i-\theta_1)^2}{2 \theta_2^2} \\
    &\overset{\mathrm{set}}{=} 0
\end{aligned}$$ Solving for $\theta_1, \theta_2$ yields $$\begin{aligned}
    \hat \theta_1 &= \bar{x} \\
    \hat \theta_2 &= \frac{\sum_{i=1}^n (x_i-\bar{x})^2}{n}
\end{aligned}$$ Plugging the numbers in yields $$\begin{aligned}
    \hat \theta_1 &= 33.4267 \\
    \hat \theta_2 &= 5.0980
\end{aligned}$$}
% \newpage
% \item Let $X_1,\dots,X_n$ be a random sample from distributions with the given probability density functions.  In each case,
% \begin{itemize}
%     \item find the maximum likelihood estimator $\hat{\theta}$,
%     \item find the method of moment estimator $\widetilde{\theta}$.
% \end{itemize}
% \begin{enumerate}
% \item $$f(x;\theta)=\dfrac{1}{\theta^2} x e^{-x/\theta}, \quad 0<x<\infty, \quad 0<\theta<\infty. $$
% \item[] Hint: Gamma distribution \\ \textcolor{red}{We start by computing the likelihood. $$\begin{aligned}
%     \mathcal{L}(\theta) &= f_{X_1, ..., X_n}(x_1, ..., x_n; \theta) \\
%     &= \prod_{i=1}^n f(x_i; \theta) \\
%     &= \prod_{i=1}^n \frac{1}{\theta^2} x_i e^{-x_i/\theta} \\
%     &= \frac{1}{\theta^{2n}} \left(\prod_{i=1}^n x_i \right) e^{-\sum_{i=1}^nx_i/\theta} 
% \end{aligned}$$ Taking the log likelihood to simplify computation, we get $$\begin{aligned}
%     l(\theta) &= \log \mathcal{L}(\theta) \\
%     &= -2n\log(\theta) + \sum_{i=1}^n \log(x_i) - \sum_{i=1}^n x_i/\theta
% \end{aligned}$$ We maximize the log likelihood by taking the first order condition $$\begin{aligned}
%     l'(\theta) &= -\frac{2n}{\theta} + \frac{\sum_{i=1}^n x_i}{\theta^2} \\
%     & \overset{\mathrm{set}}{=} 0
% \end{aligned}$$ Solving for $\theta$ yields $$\begin{aligned}
%     \hat \theta = \frac{1}{2}\bar{x}
% \end{aligned}$$Notice that this is the pdf of $\mathrm{Gamma}(2, \theta)$. By the properties of the Gamma distribution, we have $\mathbb{E}(X)= 2\theta$. Setting the first sample moment to the first population moment, we get $$\begin{aligned}
%     \bar{x} = 2\theta
% \end{aligned}$$ Solving for $\theta$ yields $$\begin{aligned}
%     \tilde \theta = \frac{\bar{x}}{2}
% \end{aligned}$$}
% \item $$f(x;\theta)=\dfrac{1}{2\theta^3} x^2 e^{-x/\theta}, \quad 0<x<\infty, \quad 0<\theta<\infty. $$
% \item[] Hint: Gamma distribution \\ \textcolor{red}{We start by computing the likelihood $$\begin{aligned}
%     \mathcal{L}(\theta) &= f_{X_1,...,X_n}(x_1, ..., x_n; \theta) \\
%     &= \prod_{i=1}^n f(x_i; \theta) \\
%     &= \prod_{i=1}^n \frac{1}{2\theta^3} x_i^2 e^{-x_i/\theta} \\
%     &= \frac{1}{2^n \theta^{3n}} \left( \prod_{i=1}^n x_i^2 \right) e^{-\sum_{i=1}^n x_i/\theta}
% \end{aligned}$$ Take the log likelihood, we get $$\begin{aligned}
%     l(\theta) &= -n \log(2) - 3n \log(\theta) + 2 \sum_{i=1}^n\log(x_i) - \sum_{i-1}^n x_i / \theta
% \end{aligned}$$ We maximize the log likelihood by taking the first order condition $$\begin{aligned} 
%     l'(\theta) &= -\frac{3n}{\theta} + \frac{\sum_{i=1}^n x_i}{\theta^2} \\
%     & \overset{\mathrm{set}}{=} 0
% \end{aligned}$$ Solving for $\theta$ yields $$\begin{aligned}
%     \hat \theta = \frac{\bar{x}}{3}
% \end{aligned}$$ Now we find the method of moments estimator. Notice that the pdf given is the pdf of $\mathrm{Gamma}(3, \theta)$. So we have $\mathbb{E}(X) = 3\theta$. By setting the first sample moment equal to the first population moment, we get $$\begin{aligned}
%     \bar{x} = 3\theta
% \end{aligned}$$ Solving for $\theta$ yields $$\begin{aligned}
%     \tilde \theta = \frac{\bar{x}}{3}
% \end{aligned}$$}
% \item $$f(x;\theta)=\dfrac{1}{2}  e^{-|x-\theta|}, \quad -\infty<x<\infty, \quad -\infty<\theta<\infty. $$
% \item[] Hint: Think about median \\ \textcolor{red}{We start by computing the likelihood $$\begin{aligned}
%     \mathcal{L}(\theta) &= f_{X_1, ..., X_n}(x_1, ..., x_n; \theta) \\
%     &= \prod_{i=1}^n f(x_i; \theta) \\
%     &= \prod_{i=1}^n \frac{1}{2} e^{-|x_i-\theta|} \\
%     &= \frac{1}{2^n} e^{-\sum_{i=1}^n |x_i-\theta|}
% \end{aligned}$$ Now notice that $\mathcal{L}$ is a decreasing function in $\sum_{i=1}^n |x_i - \theta|$. To maximize $\mathcal{L}$, we need to minimize $\sum_{i=1}^n |x_i-\theta|$. Notice that $\theta = \mathrm{median}(x_1, ..., x_n)$ minimizes $\sum_{i=1}^n |x_i-\theta|$. So the maximum likelihood estimator is $$\begin{aligned}
%     \hat \theta = \mathrm{median}(x_1, ..., x_n)
% \end{aligned}$$ Now since the distribution is symmetric around $\theta$, we know that $\mathbb{E}(X) = \theta$. Setting the first sample moment equal to the first population moment yields $$\tilde \theta = \bar{x}$$}
% \end{enumerate}
% \newpage
% \item Let $X_1,\dots,X_n$ be a random sample from distributions with the given probability density functions. In each case,
% \begin{itemize}
%     \item find the maximum likelihood estimator $\hat{\theta}$,
%     \item find the method of moment estimator $\widetilde{\theta}$ %, and
%  %   \item identify whether they are biased / unbiased estimator(s). 
% \end{itemize}


% \begin{enumerate}
%     \item $f(x;\theta) = e^{-(x-\theta)}, \quad \theta\leq x <\infty, -\infty < \theta <\infty$, zero elsewhere.  \\ \textcolor{red}{We start by computing the likelihood $$\begin{aligned}
%         \mathcal{L}(\theta) &= f_{X_1, ..., X_n}(x_1, ..., x_n; \theta) \\
%         &= \prod_{i=1}^n f(x_i; \theta) \\
%         &= \prod_{i=1}^n e^{-(x_i-\theta)} \mathbb{I}_{\{x_i\geq \theta\}} \\
%         &= e^{n\theta - \sum_{i=1}^n x_i} \mathbb{I}_{\{x_{(1) }\geq \theta\}}
%     \end{aligned}$$ Notice that $\mathcal{L}$ is increasing in $\theta$ until $\theta = x_{(1)}$, when the likelihood becomes zero. So to maximize the likelihood, we want $\theta$ to be as large as possible without the likelihood becoming $0$. Therefore, $\hat \theta = x_{(1)}$. Now we compute the first population moment $$\begin{aligned}
%         \mathbb{E}(X) &= \int_{\theta}^{\infty} xe^{-(x-\theta)} dx \\
%         &= \int_0^{\infty} (u + \theta) e^{-u} du \\
%         &= \int_0^{\infty} ue^{-u} du + \theta \int_0^{\infty} e^{-u}du \\
%         &= 1 + \theta
%     \end{aligned}$$ Setting the first sample moment equal to the first population moment, we get $$\begin{aligned}
%         \bar{x} = 1+\theta
%     \end{aligned}$$ Solving for $\theta$, we get $$\tilde{\theta} = \bar{x} - 1$$}
%     \item $f(x;\theta) = \dfrac{2x}{\theta^2}, \quad 0<x\leq \theta$, zero elsewhere.  $\theta\neq 0$. \\ \textcolor{red}{We start by computing the likelihood $$\begin{aligned}
%         \mathcal{L}(\theta) &= f_{X_1, ..., X_n}(x_1, ..., x_n; \theta) \\
%         &= \prod_{i=1}^n f(x_i; \theta) \\
%         &= \prod_{i=1}^n \frac{2x_i}{\theta^2} \mathbb{I}_{\{0 < x_i \leq \theta\}} \\
%         &= \frac{2^n \prod_{i=1}^n x_i}{\theta^{2n}} \mathbb{I}_{\{x_{(1)} > 0\}} \mathbb{I}_{\{x_{(n)} \leq \theta\}}
%     \end{aligned}$$ Notice that $\mathcal{L}$ is decreasing in $\theta$ until $\theta = x_{(n)}$, when the likelihood turns zero. So we want to minimize $\theta$ without turning the likelihood to zero. To do that, we set $$\hat \theta = x_{(n)}$$ Now we compute the first population moment. $$\begin{aligned}
%         \mathbb{E}(X) &= \int_{0}^{\theta} x \cdot \frac{2x}{\theta^2} dx \\
%         &= \frac{2}{3}\theta
%     \end{aligned}$$ Setting the first sample moment equal to the first population moment yields $$\tilde{\theta} = \frac{3\bar{x}}{2}$$}
%     \item $f(x;\theta) = \dfrac{1}{\theta} x^{(1-\theta)/\theta}, \quad 0<x<1, \quad\text{zero elsewhere}, 0<\theta<\infty$. \\ \textcolor{red}{We start by computing the likelihood $$\begin{aligned}
%         \mathcal{L}(\theta) &= f_{X_1, \dots, X_n}(x_1, \dots, x_n; \theta) \\
%         &= \prod_{i=1}^n f(x_i; \theta) \\
%         &= \prod_{i=1}^n \frac{1}{\theta} x_i^{(1-\theta)/\theta} \\
%         &= \frac{1}{\theta^n} \prod x_i^{(1-\theta)/\theta}
%     \end{aligned}$$ Now we take the log likelihood $$\begin{aligned}
%         l(\theta) &= \log \mathcal{L}(\theta) \\
%         &= -n\log(\theta) + \frac{1-\theta}{\theta} \sum_{i=1}^n \log(x_i) \\
%         &= -n\log(\theta) + \frac{\sum_{i=1}^n \log(x_i)}{\theta} - \sum_{i=1}^n \log(x_i)
%     \end{aligned}$$ Now we maximize the log likelihood by taking the first order condition $$\begin{aligned}
%         l'(\theta) &= -\frac{n}{\theta} - \frac{\sum_{i=1}^n \log(x_i)}{\theta^2} \\
%         & \overset{\mathrm{set}}{=} 0
%     \end{aligned}$$ Solving for $\theta$ yields $$\hat \theta = -\frac{\sum_{i=1}^n \log(x_i)}{n}$$ Now we find the first population moment $$\begin{aligned}
%         \mathbb{E}(X) &= \int_0^1 x \cdot \frac{1}{\theta} x^{(1-\theta)/\theta}dx \\
%         &= \frac{1}{1+\theta}
%     \end{aligned}$$ Setting the first sample moment equal to the first population moment yields $$\begin{aligned}
%         \bar{x} = \frac{1}{1+\theta} 
%     \end{aligned}$$ Solving for $\theta$ yields $$\tilde{\theta} = \frac{1-\bar{x}}{\bar{x}}$$}
%     \item $f(x;\theta) = \dfrac{1}{2},\quad  \theta-1<x<\theta+1$, $-\infty<\theta<\infty$ \\ \textcolor{red}{We start by computing the likelihood $$\begin{aligned}
%         \mathcal{L}(\theta) &= f_{X_1, ..., X_n}(x_1, ..., x_n; \theta) \\
%         &= \prod_{i=1}^n f(x_i; \theta) \\
%         &= \prod_{i=1}^n \frac{1}{2} \mathbb{I}_{\{\theta - 1 < x_i < \theta + 1\}} \\
%         &= \frac{1}{2^n} \mathbb{I}_{\{x_{(1)} > \theta - 1\}} \mathbb{I}_{\{x_{(n)} < \theta + 1\}} \\
%     \end{aligned}$$Notice that the likelihood is either $1/2^n$ or $0$. To maximize the likelihood, we need to make sure $\mathcal{L}(\theta) = 1/2^n$. So can choose $\hat\theta$ to be any number such that $$x_{(n)}-1 < \hat \theta < x_{(1)} + 1$$ Now we compute the first population moment. Since the pdf is the pdf of $\mathrm{Unif}(\theta - 1, \theta + 1)$, the first population moment is $$\begin{aligned}
%         \mathbb{E}(X) = \theta 
%     \end{aligned}$$ Setting the first sample moment equal to the first population moment yields $$\tilde{\theta} = \bar{x}$$}
% \end{enumerate}
% \newpage
% \item Let $X_1,\dots,X_n$ be a random sample from distributions with the given probability density functions. In each case, find the maximum likelihood estimator $\hat{\theta}$ and the method of moment estimator $\widetilde{\theta}$.
% \begin{enumerate}
% \item $$f(x;\theta)=\dfrac{1}{\theta^2} x e^{-x/\theta}, \quad 0<x<\infty, \quad 0<\theta<\infty. $$
% \item[] Hint: Gamma distribution
% \item $$f(x;\theta)=\dfrac{1}{2\theta^3} x^2 e^{-x/\theta}, \quad 0<x<\infty, \quad 0<\theta<\infty. $$
% \item[] Hint: Gamma distribution
% \item $$f(x;\theta)=\dfrac{1}{2}  e^{-|x-\theta|}, \quad -\infty<x<\infty, \quad -\infty<\theta<\infty. $$
% \item[] Hint: Think about median
% \end{enumerate}
%\item Let $X$ equal the weight in grams of a 52-gram snack pack of candies. Assume that the distribution of $X$ is following $\mathcal{N}(\mu, 8)$. A random sample of $n=10$ observations of $X$ yielded the following data:
%%
%\begin{center}
%\begin{tabular}{c c c c c} 
%55.95 &56.54 &57.58 &55.13& 57.48\\
%56.06 &59.93 &58.30 &52.57 &58.46
%\end{tabular}
%\end{center}
%%
%Find a method of moment estimator for $\mu$.
\end{enumerate}


	


%\thispagestyle{empty}
\end{document}
